\uuid{RnKH}
\exo7id{1508}
\titre{exo7 1508}
\auteur{barraud}
\organisation{exo7}
\datecreate{2003-09-01}
\isIndication{false}
\isCorrection{false}
\chapitre{Espace euclidien, espace normé}
\sousChapitre{Orthonormalisation}
\cpge{mpsi-mp2i.espaces-prehilbertiens.a, pcsi.espaces-prehilbertiens.a, ptsi.geometrie-plan.b, mpsi-mp2i.espaces-prehilbertiens.c, pcsi.espaces-prehilbertiens.c, mpsi-mp2i.espaces-prehilbertiens.d, pcsi.espaces-prehilbertiens.d}
\module{Algèbre}
\niveau{L2}
\difficulte{}

\contenu{
\texte{
A deux polynômes $P$ et $Q$ de $\R_{n}[X]$, on associe le nombre
  $$
  \phi(P,Q)=\int_{0}^{1}P'(t)Q'(t)dt+P(0)Q(0)
  $$
}
\begin{enumerate}
    \item \question{Montrer que $\phi$ est un produit scalaire sur $\R_{n}[X]$.}
    \item \question{Lorsque $n=2$, donner une base orthonormée pour ce produit scalaire.}
\end{enumerate}
}
